\chapter{Time Series: A First Course}\label{ch:b2:time-series}

Most macroeconomic data are time series: one country, observed quarter after quarter. Each observation resembles the one
before it, the samples are short, and many series wander without returning to a fixed level. The regression methods of
the last four chapters assumed independent observations; this chapter shows what changes when they are not. The chart
shows the main danger: two series that have nothing to do with each other.

\begin{openingchart}{Two independent simulated random walks of 200 steps (each step a standard normal draw), and the
  regression of one on the other: $R^2 = 0.548$, $t = -15.5$.}{Simulated (seeded); our computation}
\begin{tikzpicture}
\begin{axis}[omrecord, width=0.66\linewidth, height=0.38\linewidth, xmin=1, xmax=200, xlabel={Step}, ylabel={Level},
    legend style={at={(1.03,0.5)}, anchor=west, font=\scriptsize}]
  \addplot[omDef, thick] table[x=step, y=a] {figdata/bachelor-2/34-time-series/out/walks.csv};
  \addlegendentry{walk A}
  \addplot[omThm, thick, dashed] table[x=step, y=b] {figdata/bachelor-2/34-time-series/out/walks.csv};
  \addlegendentry{walk B}
\end{axis}
\end{tikzpicture}
\end{openingchart}

% figdata: walks_fit (r2 0.548, t -15.5, slope -0.552)
Each walk adds an independent random step to its previous value; the two were drawn separately. Yet the regression of B on
A has an $R^2$ of 0.55 and a $t$ statistic of $-15.5$. The pair was picked to make the point; the chapter's first listing
measures how often such a result appears by chance.

\begin{recall}
From the Year 1 volume: logarithms and growth rates, the OLS estimator, $R^2$, tests and confidence intervals. From
\cref{ch:b2:regression-inference}: robust standard errors and the size of a test. From \cref{ch:b2:business-cycles}:
detrending and the \omterm{def:b2:business-cycles:cyclical-component}{cyclical component}.
\end{recall}

\section{Trend and season}\label{sec:b2:time-series:trend}

A series $y_t$, $t = 1, \dots, T$, often has a \emph{trend}, a slow, persistent movement such as the growth of real GDP,
and a \emph{season}, a pattern that repeats each year, such as retail sales that peak in December. Official statistics
usually publish seasonally adjusted series, from which the season has been estimated and removed; the series used below
are all seasonally adjusted. A trend is handled by taking logarithms and differences: $100\Delta\ln y_t$ is the growth rate
in per cent, and four times it, for a quarterly series, the annualised rate. The deviation from a smooth trend was the
business cycle of \cref{ch:b2:business-cycles}.

\section{Autocorrelation}\label{sec:b2:time-series:autocorrelation}

\begin{definition}[Autocorrelation function]\label{def:b2:time-series:acf}
The \emph{autocorrelation function}\index{autocorrelation function} of a series gives, for each lag $k$, the correlation
$\rho_k = \operatorname{Cov}(y_t, y_{t-k})/\operatorname{Var}(y_t)$ between the series and itself $k$ periods earlier. Its
sample version is $r_k = \sum_{t=k+1}^T(y_t - \bar y)(y_{t-k} - \bar y)/\sum_{t=1}^T(y_t - \bar y)^2$.
\end{definition}

\begin{omchart}{Sample autocorrelations of US quarterly inflation (annualised) and of the unemployment rate, lags 1 to 12
  quarters, 1959--2009. Dashed lines: $\pm1.96/\sqrt{T}$, the 95 per cent band for a series with no
  autocorrelation.}{statsmodels macrodata (FRED and Bureau of Labor Statistics); our computation}
\begin{tikzpicture}
\begin{axis}[omrecord, width=0.66\linewidth, height=0.38\linewidth, ybar, bar width=4pt, xmin=0.3, xmax=12.7, ymin=-0.2,
    ymax=1.05, xtick={1,...,12}, xlabel={Lag (quarters)}, ylabel={Autocorrelation},
    legend style={at={(1.03,0.5)}, anchor=west, font=\scriptsize}]
  \addplot[fill=omDef!70, draw=omDef] table[x=lag, y=infl] {figdata/bachelor-2/34-time-series/out/acf.csv};
  \addlegendentry{inflation}
  \addplot[fill=omThm!60, draw=omThm] table[x=lag, y=unemp] {figdata/bachelor-2/34-time-series/out/acf.csv};
  \addlegendentry{unemployment}
  \addplot[black!60, dashed, sharp plot, forget plot, domain=0.3:12.7, samples=2] {1.96/sqrt(203)};
  \addplot[black!60, dashed, sharp plot, forget plot, domain=0.3:12.7, samples=2] {-1.96/sqrt(203)};
\end{axis}
\end{tikzpicture}
\end{omchart}

\omcode{code/bachelor-2/34-time-series/listings/acf.py}{1}{14}{Sample autocorrelations, with lag 1 also computed by
hand.}\label{lst:b2:time-series:acf}

\omoutput{code/bachelor-2/34-time-series/outputs/acf.txt}{Both series remember: unemployment almost perfectly from one
quarter to the next.}\label{out:b2:time-series:acf}

% code: listings/acf.py; figdata: acf (infl lag 1 0.642, lag 12 0.238; unemp lag 1 0.956, lag 12 0.260)
Inflation has a first autocorrelation of 0.64, unemployment of 0.96, and both stay above the band for twelve quarters. A
regression on such data has far fewer independent observations than its 203 quarters suggest.

\section{Stationarity}\label{sec:b2:time-series:stationarity}

\begin{definition}[Weak stationarity]\label{def:b2:time-series:stationarity}
A series is \emph{weakly stationary}\index{weak stationarity} if its mean $E[y_t]$ and variance $\operatorname{Var}(y_t)$
do not depend on $t$, and the covariance $\operatorname{Cov}(y_t, y_{t-k})$ depends only on the lag $k$.
\end{definition}

\begin{definition}[White noise]\label{def:b2:time-series:white-noise}
A \emph{white noise}\index{white noise} is a series $\varepsilon_t$ with mean zero, constant variance $\sigma^2$ and no
autocorrelation at any lag.
\end{definition}

A stationary series returns to its mean; its past is informative about its future, but less and less so at longer
horizons. A series with a trend is not stationary, nor is a random walk (\cref{sec:b2:time-series:unit-roots}). Least
squares, its standard errors and its tests keep their usual large-sample behaviour when the series involved are
stationary and their autocorrelations die out fast enough; they lose it when the series are not.

\section{Autoregressive and moving-average models}\label{sec:b2:time-series:arma}

\begin{definition}[Autoregressive process]\label{def:b2:time-series:ar}
An \emph{autoregressive process}\index{autoregressive process} of order $p$, AR($p$), is $y_t = c + \phi_1y_{t-1} + \dots
+ \phi_py_{t-p} + \varepsilon_t$ with $\varepsilon_t$ \omterm{def:b2:time-series:white-noise}{white noise}.
\end{definition}

\begin{definition}[Moving-average process]\label{def:b2:time-series:ma}
A \emph{moving-average process}\index{moving-average process} of order $q$, MA($q$), is $y_t = \mu + \varepsilon_t +
\theta_1\varepsilon_{t-1} + \dots + \theta_q\varepsilon_{t-q}$ with $\varepsilon_t$ \omterm{def:b2:time-series:white-noise}{white noise}.
\end{definition}

\begin{proposition}[Moments of an AR(1)]\label{prop:b2:time-series:ar1}
If $y_t = c + \phi y_{t-1} + \varepsilon_t$ with $|\phi| < 1$ and $y_t$ stationary, then
\[
E[y_t] = \frac{c}{1 - \phi},\qquad \operatorname{Var}(y_t) = \frac{\sigma^2}{1 - \phi^2},\qquad \rho_k = \phi^k.
\]
\end{proposition}

\begin{proof}
Take expectations: $\mu = c + \phi\mu$, so $\mu = c/(1 - \phi)$. Take variances, using that $\varepsilon_t$ is
uncorrelated with $y_{t-1}$ (which depends only on earlier shocks): $\gamma_0 = \phi^2\gamma_0 + \sigma^2$. Multiply $y_t
- \mu = \phi(y_{t-1} - \mu) + \varepsilon_t$ by $y_{t-k} - \mu$ and take expectations: $\gamma_k = \phi\gamma_{k-1}$ for $k
\ge 1$, so $\gamma_k = \phi^k\gamma_0$ and $\rho_k = \phi^k$.
\end{proof}

A shock to an AR(1) decays geometrically: after $h$ periods a fraction $\phi^h$ remains, and half of it is gone after
$\ln 0.5/\ln\phi$ periods, the shock's \emph{half-life}. An MA(1) remembers one period only: $\rho_1 = \theta/(1 +
\theta^2)$ and $\rho_k = 0$ for $k \ge 2$ (Exercise 10). An AR model is fitted by least squares of $y_t$ on its own lags.

\omcode{code/bachelor-2/34-time-series/listings/ar_fit.py}{1}{23}{AR(1) and AR(2) for US unemployment, by least squares and by
statsmodels' AutoReg.}\label{lst:b2:time-series:ar-fit}

\omoutput{code/bachelor-2/34-time-series/outputs/ar_fit.txt}{The two routes agree. The AR(1) coefficient is close to one;
the AR(2) has a large negative second lag.}\label{out:b2:time-series:ar-fit}

% code: listings/ar_fit.py
The AR(1) coefficient of 0.988 implies a half-life of 58 quarters, more than fourteen years, which is not a credible
description of how unemployment recovers from a recession. The AR(2) tells a different story: a rise in unemployment is
followed by further rises ($\phi_1 = 1.64$), then a fall ($\phi_2 = -0.68$, $t = -12.9$), the hump of a recession and its
recovery. A coefficient this close to one also raises the question of the next section.

\section{Unit roots and spurious regression}\label{sec:b2:time-series:unit-roots}

\begin{definition}[Unit root]\label{def:b2:time-series:unit-root}
A series has a \emph{unit root}\index{unit root} if it is an autoregression whose coefficients sum to one; the simplest
case is the random walk $y_t = y_{t-1} + \varepsilon_t$. Its shocks never decay, and its variance grows with $t$.
\end{definition}

\begin{definition}[Spurious regression]\label{def:b2:time-series:spurious}
A \emph{spurious regression}\index{spurious regression} is a regression between series with \omterm{def:b2:time-series:unit-root}{unit roots} that shows a large
$R^2$ and significant $t$ statistics although the series are unrelated.
\end{definition}

\omcode{code/bachelor-2/34-time-series/listings/random_walks.py}{1}{21}{Two thousand pairs of independent random walks,
regressed in levels and in differences.}\label{lst:b2:time-series:random-walks}

\omoutput{code/bachelor-2/34-time-series/outputs/random_walks.txt}{In levels, a test of size 5 per cent rejects 84 per cent
of the time; in differences, 6 per cent.}\label{out:b2:time-series:random-walks}

% code: listings/random_walks.py
In levels the $t$ test rejects the true hypothesis of no relation in 84 per cent of the samples, and the median $R^2$ is
0.17. The residuals of such a regression are themselves a random walk, and the formulas for the standard errors assume
they are not. In differences, where each series is \omterm{def:b2:time-series:white-noise}{white noise}, the test rejects 5.9 per cent of the time, as it should.

\begin{definition}[Dickey--Fuller test]\label{def:b2:time-series:dickey-fuller}
The \emph{Dickey--Fuller test}\index{Dickey--Fuller test} tests the null hypothesis of a \omterm{def:b2:time-series:unit-root}{unit root} by regressing $\Delta y_t$
on a constant (and possibly a trend), $y_{t-1}$, and lags of $\Delta y_t$; the statistic is the $t$ ratio of the coefficient
on $y_{t-1}$, compared with critical values derived under the unit-root null, which are more negative than those of the
normal distribution.
\end{definition}

Under a \omterm{def:b2:time-series:unit-root}{unit root} the $t$ ratio does not follow the usual distribution: the 5 per cent critical value is about $-2.9$ with a
constant and $-3.4$ with a constant and a trend, against $-1.65$ for a one-sided normal test. The version with lags of
$\Delta y_t$ is called augmented.

\omcode{code/bachelor-2/34-time-series/listings/adf.py}{1}{16}{Augmented \omterm{def:b2:time-series:dickey-fuller}{Dickey--Fuller tests} on US log real GDP (with a
trend) and on its growth rate.}\label{lst:b2:time-series:adf}

\omoutput{code/bachelor-2/34-time-series/outputs/adf.txt}{A \omterm{def:b2:time-series:unit-root}{unit root} is not rejected in the level of log GDP, and clearly
rejected in its growth rate.}\label{out:b2:time-series:adf}

% code: listings/adf.py
Log real GDP gives $-2.38$, short of the critical value of $-3.43$: the data cannot reject a \omterm{def:b2:time-series:unit-root}{unit root} around a trend.
Its growth rate gives $-6.97$, far beyond $-2.88$. The test has low power: it often fails to reject a \omterm{def:b2:time-series:unit-root}{unit root} in series
that are stationary but persistent, and 50 years of quarterly data rarely settle the question. The practical lesson is
narrower: a regression in levels of persistent series needs a reason to be believed, and regressions in growth rates are
the safer default.

\section{Distributed lags}\label{sec:b2:time-series:distributed-lags}

\begin{definition}[Distributed-lag model]\label{def:b2:time-series:distributed-lag}
A \emph{distributed-lag model}\index{distributed-lag model} relates $y_t$ to current and past values of a regressor:
$y_t = \alpha + \beta_0x_t + \beta_1x_{t-1} + \dots + \beta_qx_{t-q} + u_t$.
\end{definition}

\begin{definition}[Dynamic multiplier]\label{def:b2:time-series:dynamic-multiplier}
In a \omterm{def:b2:time-series:distributed-lag}{distributed-lag model} the \emph{dynamic multiplier}\index{dynamic multiplier} at horizon $h$ is $\beta_h$, the effect
on $y_{t+h}$ of a one-unit change in $x_t$ alone; the \emph{long-run multiplier} is $\beta_0 + \beta_1 + \dots + \beta_q$,
the effect of a permanent one-unit change in $x$.
\end{definition}

If a rise of one point in a tax rate lowers sales by 0.2 per cent in the same quarter, 0.3 in the next and 0.1 in the one
after, the \omterm{def:b2:time-series:dynamic-multiplier}{dynamic multipliers} are $-0.2$, $-0.3$ and $-0.1$, and the long-run effect of a permanent rise is $-0.6$ per
cent. The coefficients have a causal reading only if $x$ is unrelated to the error at all leads and lags, the conditions of
\cref{ch:b2:causality-and-experiments} stated in time; and since $u_t$ is usually autocorrelated, the standard errors need
the correction of \cref{sec:b2:time-series:hac}.

\section{Forecasting and how to score it}\label{sec:b2:time-series:forecasting}

\begin{definition}[Out-of-sample forecast]\label{def:b2:time-series:oos}
An \emph{out-of-sample forecast}\index{out-of-sample forecast} of $y_{t+h}$ is made with a model fitted only on data
available at date $t$; a \emph{pseudo out-of-sample} exercise repeats this for each $t$ in a past period, refitting each
time, and compares the forecasts with what happened.
\end{definition}

\begin{definition}[Root mean squared forecast error]\label{def:b2:time-series:rmse}
The \emph{root mean squared forecast error}\index{root mean squared forecast error} of forecasts $\hat y_{t+h}$ is
$\sqrt{\tfrac1m\sum_t(y_{t+h} - \hat y_{t+h})^2}$ over the $m$ forecasts made.
\end{definition}

A model's fit within its own sample flatters it: it was chosen to fit those data. The honest test is out of sample, against
a simple benchmark. For inflation the benchmark is the \emph{random walk}: inflation over the next four quarters equals
inflation over the past four.

\omcode{code/bachelor-2/34-time-series/listings/forecast.py}{1}{26}{Each quarter from 1985, a regression forecast of inflation
four quarters ahead, refitted on the data then available, against the random walk.}\label{lst:b2:time-series:forecast}

\omoutput{code/bachelor-2/34-time-series/outputs/forecast.txt}{The random walk wins.}\label{out:b2:time-series:forecast}

\begin{omchart}{US consumer-price inflation over four quarters, 1986--2009, and its forecasts made four quarters earlier:
  by a regression on past inflation, refitted each quarter, and by the random walk (past four-quarter inflation), per
  cent.}{statsmodels macrodata (FRED); our computation}
\begin{tikzpicture}
\begin{axis}[omrecord, width=0.66\linewidth, height=0.38\linewidth, xmin=1986, xmax=2009.5, ymin=-2.5, ymax=6.5,
    xtick={1986,1990,...,2006}, xticklabel style={/pgf/number format/1000 sep={}}, ylabel={Per cent},
    legend style={at={(1.03,0.5)}, anchor=west, font=\scriptsize}]
  \addplot[black, thick] table[x=t, y=actual] {figdata/bachelor-2/34-time-series/out/fcst.csv};
  \addlegendentry{actual}
  \addplot[omDef, thick, dashed] table[x=t, y=ar] {figdata/bachelor-2/34-time-series/out/fcst.csv};
  \addlegendentry{regression}
  \addplot[omThm, thick, dotted] table[x=t, y=rw] {figdata/bachelor-2/34-time-series/out/fcst.csv};
  \addlegendentry{random walk}
\end{axis}
\end{tikzpicture}
\end{omchart}

% code: listings/forecast.py; figdata: rmse (ar 1.772, rw 1.624, n 95)
Over 95 forecasts the regression misses by 1.77 points in root mean square, the random walk by 1.62. The regression
forecasts lie above the outcome on average (a mean error of $-0.72$): fitted on data that include the high inflation of the
1970s, it keeps predicting a return toward a higher mean that did not come. Both miss the fall of 2009 by several points.

\begin{record}[Inflation forecasts against a random walk]
% ledger: b2:atkeson-ohanian2001
Atkeson and Ohanian (2001) compared three sets of inflation forecasts built on the relation between inflation and
unemployment (a textbook model, two models of a type then common in research, and the forecasts of the Federal Reserve
Board) with the naive forecast that inflation over the next year equals inflation over the past year, over the previous 15
years. None was more accurate than the naive forecast, and they predicted the direction of the change in inflation no
better than a coin flip.
\end{record}

\section{What Granger causality does and does not mean}\label{sec:b2:time-series:granger}

\begin{definition}[Granger causality]\label{def:b2:time-series:granger}
A series $x$ \emph{Granger-causes}\index{Granger causality} a series $y$ if past values of $x$ help predict $y_t$ given
past values of $y$; the test is an $F$ test that the coefficients on the lags of $x$ are all zero in a regression of $y_t$
on lags of $y$ and of $x$.
\end{definition}

\omcode{code/bachelor-2/34-time-series/listings/granger.py}{1}{16}{Does money growth help predict inflation, and the
reverse? Four lags each.}\label{lst:b2:time-series:granger}

\omoutput{code/bachelor-2/34-time-series/outputs/granger.txt}{Neither test rejects at 5 per cent.}\label{out:b2:time-series:granger}

% code: listings/granger.py
With four quarterly lags, money growth does not help predict inflation at the 5 per cent level ($p = 0.076$), nor
inflation money growth ($p = 0.213$). The name promises more than the test delivers. \omterm{def:b2:time-series:granger}{Granger causality} is about
prediction, not about effects in the sense of \cref{ch:b2:causality-and-experiments}: a variable that responds to
expectations of the future, such as an asset price, can predict an outcome it does not cause, and a policy that offsets
shocks can cause an outcome while failing to predict it. The test answers a forecasting question; for that question it is
the right one.

\section{Serial correlation in regressions}\label{sec:b2:time-series:hac}

When the errors of a regression are autocorrelated, least squares remains unbiased if the regressors are unrelated to the
errors, but the classical and the \omterm{def:b2:regression-inference:robust-se}{heteroskedasticity-robust standard errors} are both wrong: they ignore that neighbouring
errors move together. Overlapping data make the problem certain: the change in inflation over the next four quarters,
measured each quarter, shares three quarters with its neighbour.

\begin{definition}[HAC standard error]\label{def:b2:time-series:hac}
A \emph{HAC standard error}\index{HAC standard error} (\omterm{def:b2:regression-inference:homoskedasticity}{heteroskedasticity} and autocorrelation consistent) replaces the
middle term of the sandwich formula by $\sum_t\hat u_t^2x_tx_t' + \sum_{j=1}^Lw_j\sum_t\hat u_t\hat u_{t-j}(x_tx_{t-j}' +
x_{t-j}x_t')$ with weights $w_j = 1 - j/(L + 1)$; this is the Newey--West estimator.
\end{definition}

\omcode{code/bachelor-2/34-time-series/listings/hac.py}{1}{20}{The change in inflation over the next four quarters against
unemployment now: classical and Newey--West standard errors.}\label{lst:b2:time-series:hac}

\omoutput{code/bachelor-2/34-time-series/outputs/hac.txt}{The residuals are strongly autocorrelated, and the honest standard
error is 60 per cent larger.}\label{out:b2:time-series:hac}

% code: listings/hac.py
The slope is $-0.29$: higher unemployment is followed by falling inflation. The residuals have a first autocorrelation of
0.78, and the Newey--West standard error, 0.151, is 1.6 times the classical 0.095. The $t$ statistic falls from $-3.0$ to
$-1.9$: significant at 5 per cent with the wrong standard error, not quite with the right one.

\begin{model}{AR(1) model}
\assumes $y_t = c + \phi y_{t-1} + \varepsilon_t$ with $|\phi| < 1$ and white-noise shocks: the series returns to a fixed
mean, and a shock decays by the factor $\phi$ each period.
\predicts Autocorrelations $\phi^k$; a forecast that moves from the last value toward the mean $c/(1 - \phi)$; forecast
intervals that widen with the horizon toward the series' own spread.
\tested US unemployment fits an AR(1) with $\phi = 0.988$, a half-life of 58 quarters at odds with its recoveries; an AR(2)
fits its hump far better. Inflation forecasts from regressions on past inflation did no better than a random walk from
1986 to 2009, as Atkeson and Ohanian (2001) found for richer models.
\breaks When the mean shifts (inflation in the 1970s and after 1985) or the series has a \omterm{def:b2:time-series:unit-root}{unit root}, the fitted $\phi$ and
the implied mean describe the sample, not the future; short samples cannot tell a large $\phi$ from a \omterm{def:b2:time-series:unit-root}{unit root}.
\end{model}

\section{Exercises}\label{sec:b2:time-series:exercises}

\begin{exercise}[$\star$]\label{exo:b2:time-series:1}
Read the opening chart. If you knew nothing of how the walks were made, what would the regression suggest? Why is the
suggestion wrong?
\end{exercise}

\begin{exercise}[$\star$]\label{exo:b2:time-series:2}
From the ACF chart, give the first and the twelfth autocorrelations of inflation and of unemployment, and the band. Which
series is more persistent?
\end{exercise}

\begin{exercise}[$\star$]\label{exo:b2:time-series:3}
A \omterm{def:b2:time-series:distributed-lag}{distributed-lag model} has $\beta_0 = 0.4$, $\beta_1 = 0.3$, $\beta_2 = 0.2$, $\beta_3 = 0.1$. Give the \omterm{def:b2:time-series:dynamic-multiplier}{dynamic
multipliers} and the long-run multiplier.
\end{exercise}

\begin{exercise}[$\star$]\label{exo:b2:time-series:4}
A share price predicts next quarter's GDP. Does it Granger-cause GDP? Does it cause GDP?
\end{exercise}

\begin{exercise}[$\star\star$]\label{exo:b2:time-series:5}
For an AR(1) with $\phi = 0.5$, $c = 1$, $\sigma = 1$: give the mean, the variance, the first three autocorrelations and
the half-life of a shock.
\end{exercise}

\begin{exercise}[$\star\star$]\label{exo:b2:time-series:6}
From the AR(1) output for unemployment, compute the implied mean and the half-life. Compare the implied mean with the
sample mean of the unemployment rate (5.9 per cent). What explains the gap?
\end{exercise}

\begin{exercise}[$\star\star$]\label{exo:b2:time-series:7}
\textbf{(D)} Add the unemployment rate at the forecast date to the forecasting regression of the listing and repeat the
exercise. Does the root mean squared error fall?
\end{exercise}

\begin{exercise}[$\star\star$]\label{exo:b2:time-series:8}
Why does the regression in differences of the random-walk listing reject at the nominal rate while the regression in
levels does not?
\end{exercise}

\begin{exercise}[$\star\star$]\label{exo:b2:time-series:9}
From the HAC output, compute the $t$ statistic of the slope with each standard error. Which conclusion changes at the
5 per cent level, and which does not?
\end{exercise}

\begin{exercise}[$\star\star\star$]\label{exo:b2:time-series:10}
For the MA(1) $y_t = \varepsilon_t + \theta\varepsilon_{t-1}$, derive the variance and the autocorrelations. Which
value of $\theta$ gives the largest first autocorrelation?
\end{exercise}

\begin{exercise}[$\star\star\star$]\label{exo:b2:time-series:11}
For a random walk starting at $y_0 = 0$, show that $\operatorname{Var}(y_t) = t\sigma^2$ and that the correlation of $y_t$
with $y_{t-k}$ is $\sqrt{(t - k)/t}$. Why does this rule out stationarity?
\end{exercise}

\begin{exercise}[$\star\star\star$]\label{exo:b2:time-series:12}
For an AR(1) with known $c$, $\phi$ and $\sigma$, derive the $h$-step forecast $E[y_{t+h} \mid y_t]$ and its error variance.
What do they tend to as $h$ grows?
\end{exercise}

\section{Problem: Brenland's inflation}\label{sec:b2:time-series:problem}

\begin{problem}[{Weekend problem --- Brenland's inflation: growth rates and autocorrelation, an AR(1), a unit-root test,
and a two-year forecast}]
\label{pb:b2:time-series:1}
Brenland's consumer price index for 2000--2025 is in the country's national accounts (\texttt{accounts}, column
\texttt{cpi}, 2015 = 100). Use annual inflation $\pi_t = 100\ln(\mathit{CPI}_t/\mathit{CPI}_{t-1})$, 2001--2025.

\medskip
\noindent\textbf{Part I --- Growth rates and autocorrelation.}
\begin{enumerate}
  \item How many observations of $\pi_t$ are there?
  \item Give the mean and standard deviation of inflation.
  \item In which year was inflation highest, and how high?
  \item Give inflation in 2025.
  \item Give the first sample autocorrelation.
  \item Is it outside the 95 per cent band for \omterm{def:b2:time-series:white-noise}{white noise}?
\end{enumerate}

\medskip
\noindent\textbf{Part II --- An AR(1).}
\begin{enumerate}[resume]
  \item Fit $\pi_t = c + \phi\pi_{t-1} + \varepsilon_t$ by least squares: give $c$ and $\phi$.
  \item Give the standard error of $\phi$ and its $t$ statistic against zero.
  \item Give the implied mean and compare it with the sample mean.
  \item Give the half-life of an inflation shock.
  \item Give the standard deviation of the residuals.
  \item What does the AR(1) say about the 2022 surge?
\end{enumerate}

\medskip
\noindent\textbf{Part III --- A unit-root test.}
\begin{enumerate}[resume]
  \item Write the Dickey--Fuller regression with a constant and no lags.
  \item Give the statistic.
  \item Give the 5 per cent critical value for this sample.
  \item Is a \omterm{def:b2:time-series:unit-root}{unit root} rejected at 5 per cent? At 10 per cent?
  \item How can $\hat\phi = 0.48$ fail to reject $\phi = 1$?
  \item What does this mean for Part IV?
\end{enumerate}

\medskip
\noindent\textbf{Part IV --- A forecast and its interval.}
\begin{enumerate}[resume]
  \item Forecast inflation for 2026.
  \item Forecast inflation for 2027.
  \item Give the standard error of each forecast, ignoring the uncertainty in $c$ and $\phi$.
  \item Give the 95 per cent interval for 2027.
  \item Why is the 2027 interval wider than the 2026 one, and what width does it tend to?
  \item What would a random-walk forecast give for 2027?
  \item State the result: the AR(1) persistence, the half-life of an inflation shock, and the two-year forecast with its
  interval.
\end{enumerate}
\end{problem}
